// Matrix Transform — mat3 for 2D affine UV transforms
//
// Educational example: build a rotation + scale matrix using mat3,
// apply it to UV coordinates, then sample the image at the transformed
// position. Also demonstrates transpose().
//
// The 3x3 matrix encodes 2D affine transforms in homogeneous coords:
//   | sx*cos  -sy*sin  tx |
//   | sx*sin   sy*cos  ty |
//   |   0        0      1 |
//
// Connect @image to any RGB input.

f$angle = 15.0;        // rotation in degrees
f$scale = 1.0;         // uniform scale factor
f$center_x = 0.5;      // rotation center X (0.5 = image center)
f$center_y = 0.5;      // rotation center Y

// ── Build rotation + scale matrix ───────────────────────────────────
float rad = $angle * PI / 180.0;
float c = cos(rad) * $scale;
float s = sin(rad) * $scale;

// Rotation-scale matrix (rows: x-axis, y-axis, translation)
mat3 rot = mat3(
     c, -s, 0.0,
     s,  c, 0.0,
    0.0, 0.0, 1.0
);

// Demonstrate transpose — for a pure rotation matrix, transpose = inverse
mat3 rot_inv = transpose(rot);

// ── Apply transform to UV ───────────────────────────────────────────
// Shift to center, apply inverse transform, shift back.
// We use the inverse so we're mapping output -> input (backward mapping).
float cu = u - $center_x;
float cv = v - $center_y;

// Multiply [cu, cv, 1] by the inverse rotation matrix
vec3 uv_hom = vec3(cu, cv, 1.0);
vec3 transformed = rot_inv * uv_hom;

float su = transformed.x + $center_x;
float sv = transformed.y + $center_y;

// ── Sample with boundary check ──────────────────────────────────────
vec3 result = vec3(0.0);
if (su >= 0.0 && su <= 1.0 && sv >= 0.0 && sv <= 1.0) {
    result = sample(@image, su, sv);
}

@OUT = result;
